\documentclass[12pt,a4paper]{exam}

\usepackage[utf8]{inputenc}
\usepackage[T1]{fontenc}
\usepackage[french]{babel}
\usepackage{amsmath,amssymb}
\usepackage{pgfplots}
\pgfplotsset{compat=1.18}
\usepackage{tikz}
\usetikzlibrary{arrows.meta}
\usepackage{geometry}
\geometry{margin=2cm}

\setlength{\parindent}{0pt}

\begin{document}

\begin{center}
\Large \textbf{- Interrogation -\\Écriture fractionnaire (401)}
\vspace{0.5cm}
\\\large \underline{\textbf{JUSTIFIER VOS RÉPONSES SUR VOTRE COPIE}}
\end{center}
\vspace{0.4cm}
\noindent\hrulefill
\vspace{0.5cm}

% --------------------------- %
% ------- EXERCICE 1 -------- %
% --------------------------- %

\textbf{Exercice 1 – Égalités de fractions (7 points)}
\vspace{5pt}

\begin{questions}
    \question[2] À l’aide d’un \textbf{produit en croix}, dire si les égalités suivantes
    sont vraies ou fausses.
    \vspace{5pt}
    \begin{parts}
        \part $\dfrac{3}{4} = \dfrac{9}{12}$
        \vspace{5pt}
        \part $\dfrac{5}{6} = \dfrac{15}{20}$
    \end{parts}

    \vspace{10pt}

    \question[2] À l’aide d’un \textbf{produit en croix}, écrire une égalité entre deux fraction
    pour les calculs suivants :
    \vspace{5pt}
    \begin{parts}
        \part $4 \times 12 = 6 \times 8$
        \vspace{5pt}
        \part $-12 \times 7 = -21 \times 4$
    \end{parts}

    \vspace{10pt}

    \question[3] \underline{Recopier et compléter sur votre copie} les fractions pour que 
    les égalités soient vraies :
    \vspace{5pt}
    \begin{parts}
        \part $\dfrac{4}{7} = \dfrac{\ldots}{21}$
        \vspace{5pt}
        \part $\dfrac{-12}{42} = \dfrac{2}{\ldots}$
        \vspace{5pt}
        \part $\dfrac{\ldots}{-18} = \dfrac{7}{3}$
    \end{parts}
\end{questions}

\noindent\hrulefill
\vspace{0.5cm}

% --------------------------- %
% ------- EXERCICE 2 -------- %
% --------------------------- %

\textbf{Exercice 2 – Encadrement (5 points)}
\vspace{5pt}

\begin{questions}
    \question[2] Ranger les nombres suivants dans l’ordre croissant :
    \[
    \dfrac{3}{4} \quad ; \quad 0,8 \quad ; \quad \dfrac{5}{6} \quad ; \quad \dfrac{-11}{4} 
    \quad ; \quad -\dfrac{7}{5}
    \]

    \question[3] Encadrer chaque nombre entre deux entiers \textbf{consécutifs} :
    \begin{parts}
        \part $\dfrac{7}{3}$
        \part $\dfrac{14}{4}$
        \part $\dfrac{-11}{2}$
    \end{parts}
\end{questions}

\noindent\hrulefill
\vspace{0.5cm}

% --------------------------- %
% ------- EXERCICE 3 -------- %
% --------------------------- %

\textbf{Exercice 3 – Additions et soustractions de fractions (8 points)}
\vspace{5pt}

Calculer et donner les résultats sous forme de \textbf{fractions simplifiées}.
\\\underline{Détailler les calculs}.

\vspace{5pt}

\begin{questions}
    \question $A = \dfrac{11}{7} + \dfrac{3}{7}$
    \vspace{5pt}
    \question $B = \dfrac{3}{4} - \dfrac{5}{16}$
    \vspace{5pt}
    \question $C = -\dfrac{2}{3} + \dfrac{4}{12}$
    \vspace{5pt}
    \question $D = \dfrac{5}{2} - \dfrac{-30}{4} + 2$
\end{questions}

\vspace{10pt}

\noindent\hrulefill
\vspace{0.5cm}

% --------------------------- %
% ------- EXERCICE 3 -------- %
% --------------------------- %

\textbf{Exercice 4 – BONUS (??? points)}
\vspace{5pt}

Après de longues négociations, il a été convenu que Léa héritera de deux quinzièmes de la 
fortune de son oncle du bout du monde. Florian héritera d'un neuvième de cette fortune.
\\Jean et Justine se partageront équitablement le reste.
\vspace{10pt}
\\
Quelles seront les parts respectives de Jean et Justine ?
\vspace{10pt}

\noindent\hrulefill
\vspace{0.5cm}


\begin{center}
\textbf{Fin de l'interrogation — relire avant de rendre.}    
\end{center}

\end{document}
